%! $n$th Roots \& Rational Exponents — Unit 5, Lesson 5.1 — Group Activity
\documentclass[10pt]{article}
\usepackage{algebra2-article}
\usepackage{algebra2-boxes}
\newcommand{\TallMath}[1]{$\displaystyle #1\rule[-1.4em]{0pt}{3.2em}$}

\begin{document}

\pageheader{Unit 5, Lesson 5.1}{Group Activity}
\namepartnerperiod

\vspace{0.08in}

\begin{headlinebox}{forestbg}
\small\textbf{One Expression, Two Costumes.} Every radical can be written as an exponent, and every
rational exponent can be written as a radical --- and \emph{which costume you choose decides how
hard the problem is}. With your partner, translate carefully and say out loud which form you picked
and why. Unless a problem says otherwise, assume all variables represent \textbf{positive} real
numbers.
\end{headlinebox}

\vspace{0.06in}

\begin{tcolorbox}[colback=white, colframe=black!40, title=\textbf{Tier R --- Remediate},
    fonttitle=\bfseries, arc=1.5mm, left=3mm, right=3mm, top=2mm, bottom=2mm, breakable]
\textbf{1. Translate.} Fill in every empty cell. The \textbf{index} always becomes the
\textbf{denominator}.
\begin{center}
\renewcommand{\arraystretch}{1.9}
\begin{tabularx}{\linewidth}{@{}X X X@{}}
\hline
\textbf{Radical form} & \textbf{Rational exponent form} & \textbf{Value (numbers only)} \\
\hline
$\sqrt{x}$            & \blank{2.4cm} & --- \\
$\sqrt[3]{x^{2}}$     & \blank{2.4cm} & --- \\
\blank{2.4cm}         & $x^{5/2}$     & --- \\
\blank{2.4cm}         & $t^{1/6}$     & --- \\
$\sqrt[4]{16}$        & \blank{2.4cm} & \blank{2.0cm} \\
\blank{2.4cm}         & $8^{1/3}$     & \blank{2.0cm} \\
$\left(\sqrt[3]{27}\right)^{2}$ & \blank{2.4cm} & \blank{2.0cm} \\
\hline
\end{tabularx}
\end{center}

\vspace{4pt}
\textbf{2. Where does each number go?} In the expression $a^{m/n}$:
\begin{itemize}[leftmargin=1.6em, itemsep=4pt]
  \item the number $n$ (the \emph{bottom}) is the \blank{3.0cm};
  \item the number $m$ (the \emph{top}) is the \blank{3.0cm}.
\end{itemize}

\vspace{2pt}
\textbf{3. Evaluate.} Take the \emph{root} first, then the power --- the numbers stay small.
\begin{center}
\renewcommand{\arraystretch}{1.7}
\begin{tabular}{@{}l l l@{}}
(a) \; $9^{1/2}=$ \blank{1.8cm} & (b) \; $64^{1/3}=$ \blank{1.8cm} &
(c) \; $4^{3/2}=$ \blank{1.8cm} \\
(d) \; $25^{3/2}=$ \blank{1.8cm} & (e) \; $8^{4/3}=$ \blank{1.8cm} &
(f) \; $(-27)^{1/3}=$ \blank{1.8cm} \\
\end{tabular}
\end{center}

\vspace{2pt}
\textbf{4.} One of the six above had a \emph{negative} base and still worked. Which one, and what is
it about its \textbf{denominator} that allowed it? \writeline
\end{tcolorbox}

\begin{tcolorbox}[colback=white, colframe=black!40, title=\textbf{Tier A --- Approaching Proficiency},
    fonttitle=\bfseries, arc=1.5mm, left=3mm, right=3mm, top=2mm, bottom=2mm, breakable]
\textbf{1. No calculator.} Evaluate each. Show the root you took and the power you applied.
\begin{center}
\renewcommand{\arraystretch}{2.0}
\begin{tabularx}{\linewidth}{@{}X X X@{}}
(a) \; $27^{4/3}=$ \blank{2.2cm} &
(b) \; $16^{3/4}=$ \blank{2.2cm} &
(c) \; $(-8)^{2/3}=$ \blank{2.2cm} \\
(d) \; $81^{-1/2}=$ \blank{2.2cm} &
(e) \; $32^{-3/5}=$ \blank{2.2cm} &
(f) \; $\left(\dfrac{8}{27}\right)^{2/3}=$ \blank{2.2cm} \\
\end{tabularx}
\end{center}

\vspace{3pt}
\textbf{2. Justify --- why root first?} Your partner insists on doing (b) as $\sqrt[4]{16^{3}}$.
That is not \emph{wrong}. Carry it out far enough to see what happens, then explain why taking the
root first is the better habit. \writelines{3}

\vspace{2pt}
\textbf{3. Justify --- the sign trap.} A classmate says ``$81^{-1/2}$ must be negative, because the
exponent is negative.'' They are wrong. Explain what the negative exponent actually does, and give
the correct value. \writelines{3}

\vspace{2pt}
\textbf{4. Sort them.} Without evaluating, decide whether each expression is a \emph{real number},
and give your reason (look at the denominator of the exponent and the sign of the base).
\begin{center}
\renewcommand{\arraystretch}{1.7}
\begin{tabularx}{\linewidth}{@{}l X X@{}}
\hline
\textbf{Expression} & \textbf{Real number? (yes/no)} & \textbf{Reason} \\
\hline
$(-64)^{1/3}$ & \blank{2.0cm} & \blank{5.0cm} \\
$(-64)^{1/2}$ & \blank{2.0cm} & \blank{5.0cm} \\
$(-32)^{1/5}$ & \blank{2.0cm} & \blank{5.0cm} \\
$(-16)^{3/4}$ & \blank{2.0cm} & \blank{5.0cm} \\
\hline
\end{tabularx}
\end{center}
\end{tcolorbox}

\begin{tcolorbox}[colback=white, colframe=black!40, title=\textbf{Tier E --- Extension},
    fonttitle=\bfseries, arc=1.5mm, left=3mm, right=3mm, top=2mm, bottom=2mm, breakable]
\textbf{Part 1 --- Problems radicals cannot do.} Each product or quotient below has \emph{mismatched
indices}, so no radical rule applies. Rewrite with rational exponents, simplify, and convert back to
a single radical.
\begin{center}
\renewcommand{\arraystretch}{2.1}
\begin{tabularx}{\linewidth}{@{}X X@{}}
(a) \; $\sqrt{x}\cdot\sqrt[4]{x} = $ \blank{4.0cm} &
(b) \; $\dfrac{\sqrt[3]{x^{2}}}{\sqrt{x}} = $ \blank{4.0cm} \\
(c) \; $\left(\sqrt[3]{x^{2}}\right)^{3/2} = $ \blank{4.0cm} &
(d) \; $\sqrt{\sqrt[3]{x}} = $ \blank{4.0cm} \\
\end{tabularx}
\end{center}
For (d), explain what happened to the two indices and why that makes sense. \writelines{2}

\vspace{5pt}
\textbf{Part 2 --- A rational exponent in the wild.} Kepler's third law says that a planet's
\textbf{orbital period} $T$ (in Earth years) depends on its \textbf{average distance} $a$ from the
Sun (in astronomical units, where $1$~AU is Earth's distance):
\[
  T(a) = a^{3/2}.
\]
\begin{center}
\renewcommand{\arraystretch}{1.25}
\begin{tabular}{|c|c|c|c|c|}
\hline
$a$ (AU)        & $1$ & $4$ & $9$ & $25$ \\ \hline
$T(a)$ (years)  & $1$ & \blank{1.2cm} & \blank{1.2cm} & \blank{1.2cm} \\ \hline
\end{tabular}
\end{center}
\begin{enumerate}[label=(\alph*), leftmargin=1.9em, itemsep=6pt]
  \item Complete the table \emph{without a calculator}. Which root did you take, and which power did
        you apply? \writeline
  \item Write $T(a)=a^{3/2}$ in radical form: $T(a)=\blank{3.0cm}$. \; Which of the two
        equivalent radical forms did you use, and why is it the easier one to compute with?
        \writeline
  \item What is the domain of $T$ \emph{algebraically} (look at the denominator of the exponent)?
        \blank{3.0cm} \\
        Does the context of the problem restrict it any further? Explain. \writelines{2}
  \item Jupiter sits about $5.2$~AU from the Sun. Between which two \emph{table entries} must its
        orbital period fall, and roughly how many Earth years is that? \writeline
  \item Solve $T=a^{3/2}$ for $a$ in terms of $T$. \quad $a = \blank{3.0cm}$ \\
        \emph{Hint:} raise both sides to the power that undoes $\frac{3}{2}$. Which power is that,
        and why does it work? \writelines{2}
  \item Your answer to (e) undoes $T(a)$ --- it is the \textbf{inverse} relationship from Lesson 5.0,
        now written with exponents instead of graphs. Use it to check one entry from the table, and
        state what the exponent $\frac{2}{3}$ is doing in words. \writelines{2}
\end{enumerate}
\end{tcolorbox}

\end{document}
